%& -shell-escape
% !TeX root = thesis-abstract.tex
% !TeX encoding = UTF-8
% !TeX program = XeLaTeX

\documentclass{.local/lib/classes/ndvr-super-article-standalone}

\begin{document}
    \begin{description-item}{Актуальность темы}
        Актуальность проблемы поиска нечётких
        дубликатов видео объясняется
        интенсивным использованием видеоданных 
        в~современных вычислительных
        системах, что закономерно порождает проблему эффективного
        их~поиска и~быстрого доступа к ним.

        % За 20-30 лет человечество накопило большое 
        % количество видеоинформации.
        % Речь идет не только о видео в Интернете. 
        % Речь идет, в первую очередь, о камерах наблюдения.
        % Этот массив информации невозможно просмотреть целиком.
        % Эффективных методов анализа видео, 
        % подобных анализу текста, 
        % в настоящее время нет. 

        Начиная с середины 2000-х годов в Интернете
        стало появляться много видео с~очень похожим содержанием.
        Такие видео называют почти-дубликатами
        (от~английского термина «near duplicate»), 
        полудубликатами,
        неполными или нечёткими дубликатами.
        Мы будем придерживаться термина «нечёткие дубликаты».
        Слово «нечеткий» в данном случае 
        означает частичное совпадение
        оригинала и дубликата.

        Задача эффективной идентификации нечётких 
        дубликатов играет ключевую
        роль 
            в~задачах поиска, 
            защите авторских прав, 
            идентификации событий
            и~многих других.


        Задача поиска нечётких дубликатов успешно решается 
        для текстов и для статических изображений.
        На первый взгляд, задача поиска нечётких дубликатов видео
        может быть сведена к задаче поиска нечётких
         дубликатов изображений.

        Однако это не так.
        Видео — это не набор статических изображений, 
        а последовательность с заданным порядком. 
        Если решать задачу поиска похожих видео явным образом, 
        придется сравнивать одну последовательность 
        изображений с другой. 
        Два одинаковых набора кадров
        с разным порядком следования
        — образуют разные видео.
        Ситуация показана на рисунке
        \ref{fig:video-sequence-shuffled}.

        \begin{figure}[H]
            \includegraphics[width=\textwidth]
                {.local/vec/out/pdf/video-sequence-shuffled.pdf}
            \caption{%
                Сравнение исходного видео %
                и~видео с перемешанными кадрами
            }
            \label{fig:video-sequence-shuffled}
        \end{figure}


        На рисунке \ref{fig:video-sequence-punctured} 
        показана обратная ситуация. 
        В последовательности кадров 
        случайным образом удалили несколько кадров. 
        Порядок следования кадров при этом сохранили. 
        Исходное и полученное видео будут 
        выражать одно и то же явление. 
        Они будут нечёткими дубликатами.


        \begin{figure}[H]
            \includegraphics[width=\textwidth]
                {.local/vec/out/pdf/video-sequence-punctured.pdf}
            \caption{%
                Сравнение исходного видео %
                и~видео с удаленными кадрами
            }
            \label{fig:video-sequence-punctured}
        \end{figure}


        Обычно полнометражный фильм содержит примерно 
        140-180 тысяч кадров. 
        При наивном подходе для поиска одного 
        видео потребуется сделать 140-180 тысяч поисковых запросов. 
        Даже при условии, что на поиск каждого 
        изображения потребуется 10 миллисекунд\footnote{
            Бабенко A. В. 
            Эффективные алгоритмы поиска по большим коллекциям изображений — 
            Диссертация на соискание учёной степени кандидата 
            физико-математических наук
            Специальность 05.13.11, 
            — М.: Московский физико-технический институт, 2017.
        }, 
        на поиск видео целиком потребуется около получаса.


        Задача поиска нечётких дубликатов видео
        похожа на задачу поиска 
        наибольшей общей подпоследовательности.
        Но в задачах сравнения текстовых последовательностей 
        требуется конечный фиксированный алфавит\footnote{
            Дэн Гасфилд — Строки, деревья и последовательности в алгоритмах: 
            Информатика и вычислительная биология — СПб.: 
            Невский Диалект; БХВ-Петербург, 2003. — 654 с: ил.        
        }.
        Несмотря на то, что существует потенциально бесконечное
        множество всевозможных изображений, 
        кадры могут быть представлены в виде символов\footnote{
            Wang Y., Belkhatir M., Tahayna B. 
            Near-duplicate Video Retrieval Based on Clustering 
            by Multiple Sequence Alignment 
            // Proceedings of the 20th ACM International Conference 
            on Multimedia. New York, NY, USA: ACM, 2012. pp. 941–944.
        }:
        \begin{itemize}
            \item   кадр огрубляют до заданного числа блоков 
                    c усредненной яркостью;
            \item   каждому блоку выдается его порядковый номер 
                    по значению его яркости 
                    относительно яркостей других блоков;
            \item   кадр представляют как кортеж 
                    порядковых номеров блоков;
            \item   полученному кортежу 
                    ставится в соответствие 
                    символ из некоторого алфавита.
        \end{itemize}

        На рисунке \ref{fig:video-dna} показана схема перехода 
        от кадров видео к символам латинского алфавита.

        \begin{figure}[h]
            \includegraphics[width=0.9\textwidth]
                {.local/vec/out/pdf/video-dna.pdf}
            \caption{%
                Схема формирование символьного\ %
                представления видео%
            }
            \label{fig:video-dna}
        \end{figure}

        Переход от видео к последовательности символов 
        — искусственное допущение, которое не отражает 
        ни структуру, ни семантику данных.
        Не все кадры видео имеют одинаковую ценность. 
        Потому сравнивать все кадры смысла не имеет. 
        Нужно выбрать только те изображения, которые наиболее 
        точно характеризуют происходящее в видео.
        Для этой цели в работе предлагается использовать
        кадры со сменой сцен.
        Смены сцен в~видео не происходят случайно. 
        Чаще всего это элементы развития сюжета.
        Если изображение в видео поменялось — это означает, 
        что произошли некоторые события, 
        которые и привели к изменениям. 
        Это показано на рисунке \ref{fig:video-sequence-events}.

        \begin{figure}[H]
            \includegraphics[width=\textwidth]
                {.local/vec/out/pdf/video-sequence-events.pdf}
            \caption{%
                Начальные кадры смен сцен из начала мультфильма\\ 
                «Дядя Стёпа~—~милиционер»
            }
            \label{fig:video-sequence-events}
        \end{figure}

        В работе предложен алгоритм поиска нечётких дубликатов,
        в котором, 
            видео сначала сравниваются 
            по времени между «событиями»~— сменами сцен,
            а потом по содержимому кадров во время смены. 
        Для видео с разной последовательностью 
        событий сравнение кадров не потребуется.
        Для видео с одинаковой последовательностью событий 
        нужно будет сравнить не все кадры, а только некоторые.

        Переход от последовательности кадров 
        к последовательности событий сводит задачу индексации 
        и поиска видео к задаче индексации 
        и поиска одномерных временных рядов. 

        Такое сведение позволяет сильно ускорить 
        поиск нечётких дубликатов видео
        и уменьшить его вычислительную сложную сложность.
        Именно в этим обусловлена актуальность представленной работы.

        Кроме того, на данный момент не существует
        высокоуровневых средств обработки видео сигналов,
        которые были бы одновременно удобны 
        и в исследовательских и в промышленных задачах\footnote{
            Солонина А.И., Клионский Д.М., Меркучева Т.В., 
            Перов С.Н., Цифровая обработка сигналов и MATLAB, 2013г.
        }.
        Поэтому в процессе реализации алгоритма поиска 
        нечётких дубликатов видео был 
        разработан декларативный алгоритмический язык
        фильтрации видео.


    \end{description-item}
\end{document}
